\documentclass[11pt]{article}
\usepackage[T1]{fontenc}\usepackage{lmodern}
\usepackage[margin=1in]{geometry}
\usepackage{amsmath,amssymb,amsthm,hyperref}
\newtheorem{proposition}{Proposition}
\setlength{\parskip}{5pt}
\title{Machine Markets: Economic-Contract Repairs and Exact Counterexamples}
\author{Research revision evidence}\date{October 7, 2026}
\begin{document}\maketitle

\section{Conditional Ambiguity Pricing and Allocation Surplus}
Fix a one-unit ticket and a CARA dealer with utility $-\exp(-\gamma\pi)$, $\gamma>0$. Under a candidate model $\theta$, the ticket's monetary adverse mark is $Z_\theta\sim N(\mu_\theta,\sigma_\theta^2)$. A total spread receipt $s$ produces profit $s-Z_\theta$. All costs below are monetary per this fixed ticket; other ticket quantities require the corresponding exposure scaling. Set
\[
c_\theta=\mu_\theta+\frac{\gamma\sigma_\theta^2}{2},\qquad
c_2=\sup_{\theta\in\Theta}c_\theta<\infty.
\]
Assume $c_\theta\ge0$ and nonnegative quotes. Competition imposes zero robust certainty equivalent, with no additional fixed costs or inventory obligations in this ticket experiment. The informed benchmark knows $\theta_0\in\Theta$ and has $c_1=c_{\theta_0}$.
\begin{proposition}[Conditional robust quote and zero-premium cases]
The minimum robust break-even quote is $s_2=c_2$, while $s_1=c_1$. Thus $\Delta=s_2-s_1\ge0$, with strict inequality if and only if $\sup_\theta c_\theta>c_{\theta_0}$. Parameter nonidentification alone does not imply a positive premium, nor proportionality to a coordinate-dependent parameter diameter.
\end{proposition}
\begin{proof}
The Gaussian moment-generating function gives
\[
-\gamma^{-1}\log\mathbb E_\theta\exp[-\gamma(s-Z_\theta)]
=s-\mu_\theta-\gamma\sigma_\theta^2/2=s-c_\theta.
\]
Its infimum over models is $s-c_2$. A nonnegative robust certainty equivalent therefore requires $s\ge c_2$. The informed calculation substitutes $\theta_0$. If every observationally equivalent latent decomposition has the same relevant cost, $c_\theta=c_1$, then $\Delta=0$ despite nonidentification. Rescaling coordinates of $\theta$ changes a Euclidean parameter diameter without changing any $c_\theta$ or quote.
\end{proof}

For an additional declared allocation-surplus criterion, let inverse demand be $v(V)=(V_0-V)/b$, $b>0$, and fixed effective marginal cost be $c$. Both volumes must be interior: $0\le V_0-b(s_0+\Delta)\le V_0-bs_0$. Treating a risk charge as this cost is an extra welfare assumption; transfers and risk sharing can change the relevant social criterion.
\begin{proposition}[Allocation loss with an existing spread distortion]
The surplus lost when the quote rises from $s_0$ to $s_0+\Delta$ is
\[
\int_{V_0-b(s_0+\Delta)}^{V_0-bs_0}[v(V)-c],dV
=b(s_0-c)\Delta+\frac b2\Delta^2.
\]
The triangle $b\Delta^2/2$ requires the efficient baseline $s_0=c$. The result does not by itself identify an empirical welfare cost of a cohort-equivalence class.
\end{proposition}
\begin{proof}
Substitute $v(V)=(V_0-V)/b$ and integrate the affine integrand over an interval of length $b\Delta$. Its endpoint values are $s_0-c$ and $s_0+\Delta-c$; interval length times their average gives the formula.
\end{proof}

\section{An Explicit Capacity Choice and Its Stability Conditions}
Let $A=E-\mu|I|P$ be available monetary capital, with positive mark $P$ and inventory $I$ in quantity units. A static capacity illustration fixes net execution edge $r\ge0$, holding-cost coefficient $h>0$, capital per posted quantity $m>0$, and quantity bounds $0\le L_{\min}\le L_{\max}$. Its problem is
\[
\max_{L\in[L_{\min},L_{\max}]}\{rL-hL^2/2\},\qquad mL\le A.
\]
Here $r$ has monetary/quantity units, $h$ monetary/quantity-squared units, and $m$ monetary/quantity units. Demand, fills and inventory costs are fixed primitives of this one-period illustration; it is not the solution to an unspecified dynamic competitive-dealer problem.
\begin{proposition}[Capacity choice and quoting infeasibility]
If $\bar L=\min(L_{\max},A/m)<L_{\min}$, the quoting commitment is infeasible. Otherwise the unique optimum is
\[
L^*=\min\{\bar L,\max(L_{\min},r/h)\}.
\]
For a declared linear book spanning a price width $w>0$, slope depth is $D=L^*/w$ and impact is $\kappa=w/L^*$ when $L^*>0$. Increasing capital weakly raises capacity; the capacity is unchanged when its capital cap is nonbinding.
\end{proposition}
\begin{proof}
The objective is strictly concave and has derivative $r-hL$. Its unrestricted maximizer is $r/h$. Projection onto the nonempty feasible interval gives the unique solution. Its upper endpoint weakly increases with $A$, proving the comparative static. The price-impact identity is a declared linear clearing schedule, not inferred from the optimization alone.
\end{proof}

With fixed anchor response $\varphi>0$, the depth-dependent price coefficient is $a=1-\kappa\varphi$. For one cohort with $0<\eta<1$ and weight $q$, the transition is
\[
J(\kappa)=\begin{pmatrix}1-\kappa\varphi&\kappa q\\
\eta(1-\kappa\varphi)&1-\eta+\eta\kappa q\end{pmatrix}.
\]
\begin{proposition}[Exact single-cohort depth stability boundary]
For $\kappa>0$, $\varphi>0$ and $0<\eta<1$, $J(\kappa)$ is Schur stable precisely when
\[
q<\varphi,\qquad
\kappa<\frac{4-2\eta}{(2-\eta)\varphi-\eta q}.
\]
If the earlier restriction $0\le a<1$ is maintained, it additionally requires $0<\kappa\varphi\le1$. Higher impact can therefore move eigenvalues as well as transient gains. It does not necessarily create large non-normal amplification.
\end{proposition}
\begin{proof}
Let $T=2-\eta+\kappa(\eta q-\varphi)$ and $d=(1-\eta)(1-\kappa\varphi)$ be trace and determinant. The three strict Jury inequalities are
\[
1-T+d=\kappa\eta(\varphi-q)>0,
\]
\[
1+T+d=4-2\eta-\kappa[(2-\eta)\varphi-\eta q]>0,
\]
\[
1-d=\eta+(1-\eta)\kappa\varphi>0.
\]
The first requires $q<\varphi$. This makes the denominator in the second bound positive, while the third is automatic.
\end{proof}

An exact positive-feedback counterexample uses $\varphi=1$, $q=1/10$, $\eta=9/10$. At $\kappa=1/10$,
\[
J=\begin{pmatrix}9/10&1/100\\81/100&109/1000\end{pmatrix},
\quad \|Je_1\|_2^2=14661/10000>1.
\]
At the larger impact $\kappa=1$,
\[
J=\begin{pmatrix}0&1/10\\0&19/100\end{pmatrix},
\quad \|J\|_2^2=461/10000<1.
\]
Both are stable by the preceding inequalities. In the fixed Euclidean all-state experiment, the lower-impact model has $M\ge\sqrt{14661}/100>1$, whereas the higher-impact model has $M=1$, since every positive-time power contracts and $J^0=I$. Thus a universal depth-thinning/transient-danger implication is false even with positive feedback. This does not invalidate amplification results under their own fixed-spectrum and output assumptions.

\section{Exact Cash, Inventory and Marked Equity}
Let signed aggressive quantity $q_t>0$ mean the dealer sells, with the trade occurring at the current midpoint $P_t$ plus signed half-spread. Let fees $f_t$ and funding charges $g_t$ be monetary. With current inventory $I_t$ and cash $C_t$,
\[
I_{t+1}=I_t-q_t,\qquad
C_{t+1}=C_t+q_tP_t+s_t|q_t|-f_t-g_t.
\]
Mark equity after the subsequent price change at $P_{t+1}$.
\begin{proposition}[Self-financing marked-equity identity]
For $E_t=C_t+I_tP_t$,
\[
E_{t+1}-E_t=s_t|q_t|+(I_t-q_t)(P_{t+1}-P_t)-f_t-g_t.
\]
The old expression $s_t|q_t|-I_t(P_{t+1}-P_t)$ is incompatible with this signed trade and inventory convention.
\end{proposition}
\begin{proof}
Substitute the cash and inventory updates into $E_{t+1}=C_{t+1}+I_{t+1}P_{t+1}$ and cancel $q_tP_t$ against the inventory trade. The resulting mark exposure is the post-trade inventory. A long inventory with no trade gains when the price rises.
\end{proof}

A fixed model loss budget $B$ in monetary units is distinct from posted quantity and a norm radius in latent-state coordinates. Choosing $0<B\le A_0$ does not ensure future margin feasibility when positions or prices change. Conditional transient/input bounds require a financially specified output operator, and actual execution additionally requires fees, funding, currency conversion, margin rules and fundamental-price risk. These are calibration and information gates, not consequences of an algebraic eigenvalue test.

\end{document}
